\section{Part 11a-5: the quantitative maximal-case theorem and its budget}\label{bpP:p11a5}
\begingroup\sloppy

This is the last of the five files of Part~11a. Files~11a-1 to 11a-4 prove, one ingredient at a time, that an analytic Hamiltonian $H=h(I)+f(\theta,I)$ whose frequency map $\omega=\nabla h$ has finite type and whose perturbation $f$ is small has a large set of invariant Lagrangian tori that move little and depend Lipschitz-continuously on the pair (action, angle). This file only assembles them. It states the theorem in the exact form in which Part~11 consumes it (Definition~\ref{bpP:p11:kam-input}, items (K1)--(K6)), fixes the two numbers $\gamma$ and $\varepsilon_{\mathrm{adm}}$ by explicit formulas of the data, proves that $\varepsilon_{\mathrm{adm}}$ may be replaced by a function monotone in every datum and bounded below by a power of the data, and records the consequences Part~13 needs (Corollary~\ref{bpP:p11a:physical-interface}).

\paragraph{Order of choices.} The data come first: dimension and exponents $(N,\mu,\rho,\tau_D)$, the compact parameter set $K$ with its thickness $\vartheta$, the angular strip width $s$, the index data $(\beta,M,M_0)$, the bound $C_1$ for the Hessian, and the permitted loss of measure $\varepsilon^*$. From them alone one computes the Diophantine constant $\gamma$ and the admissible size $\varepsilon_{\mathrm{adm}}$ (Definition~\ref{bpP:p11a:thresholds}). Only then is a Hamiltonian chosen, with $\sup\lvert f\rvert$ at most $\varepsilon_{\mathrm{adm}}$. Neither the Diophantine property of the frequencies, nor the bi-Lipschitz bound, nor the smallness of the displacement is a hypothesis of anything below; all three are outputs of Files~11a-3 and 11a-4. The initial action radius $s_0$ of the iteration is a choice made inside the proof (Lemma~\ref{bpP:p11a:free-radius}), not a datum.

\paragraph{What is proved.}
\begin{itemize}
\item \texttt{thresholds} (Definition~\ref{bpP:p11a:thresholds}): the explicit $\gamma$ and the monotone threshold $\varepsilon_{\mathrm{adm}}^\flat$.
\item \texttt{budget} (Proposition~\ref{bpP:p11a:budget}): $\varepsilon^\flat_{\mathrm{adm}}\le\varepsilon_{\mathrm{adm}}$ of File~11a-3, monotonicity in every datum, and the polynomial lower bounds (general rates, and the form (K5)).
\item \texttt{zero-case} (Lemma~\ref{bpP:p11a:zero-case}): the case $f\equiv0$, which the iteration (stated for a positive bound $\varepsilon$) does not cover.
\item \texttt{weak-kam-quantitative} (Theorem~\ref{bpP:p11a:weak-kam-quantitative}): the theorem, with (K1)--(K6).
\item \texttt{physical-interface} (Corollary~\ref{bpP:p11a:physical-interface}): the set $S$ of finite positive measure, the displacement, the invariance, in the form consumed by Parts~11 and 13.
\end{itemize}

\subsection{The thresholds and the budget}\label{bpP:p11a:sec-q-budget}

\paragraph{Conventions.} Those of Files~11a-3 and 11a-4: period-one angles, max norms, $\lvert k\rvert_1$, $T_s=\{\theta\in\C^N:\lvert\operatorname{Im}\theta\rvert_\infty<s\}$ (mod $\Z^N$), $U=K+2\vartheta$. The constants $\varpi$, $C_Q$, $c_N$, $c'_{N,\mu}$ depend on $(N,\mu,\tau)$ only (Definition~\ref{bpP:p11a:schedule-consts}, Lemma~\ref{bpP:p11a:smooth-cutoff}); $A_0=\log(16c_N/\varpi)$ and $\varsigma:=s/24$ ($=\varsigma_0$ of File~11a-3). The data are
\[\mathfrak d=(N,\mu,\rho,\tau,K,\vartheta,s,\beta,M,M_0,C_1,\varepsilon^*),\]
with $\tau=\tau_D$, subject to the standing requirements of Assumption~\ref{bpP:p11a:sel-data} (a)--(d) on the numbers (not on $h,f$): $\rho\in\{1,2\}$, $\tau>N\mu/\rho-1$, $\tau\ge0$, $\vartheta,s\in(0,1]$, $C_1\ge1$, $\varepsilon^*\in(0,\operatorname{meas}K)$, $0<\beta\le M$.

\begin{definition}[$\gamma$ and the monotone threshold $\varepsilon^\flat_{\mathrm{adm}}$]\label{bpP:p11a:thresholds}
Let $d=\operatorname{diam}_\infty K$ and $s_\tau=\rho(\tau+1)/\mu-N>0$. Put
\[C_{\mathrm{Dio}}=2^{N+2}\mu(\sqrt Nd)^{N-1}\frac M\beta\Bigl(1+(4+\tfrac1\vartheta)\sqrt Nd\Bigr)\beta^{-1/\mu}N^{\rho/2\mu}\,2\cdot3^{N-1}\Bigl(1+\frac1{s_\tau}\Bigr),\]
\[\gamma_*=\min\Bigl\{\tfrac12N^{-1/2}(\beta/4)^{1/\rho},\ (\varepsilon^*/C_{\mathrm{Dio}})^{\mu/\rho}\Bigr\},\qquad \gamma=\gamma(\mathfrak d):=\tfrac12\gamma_*\]
(Corollary~\ref{bpP:p11a:power-law}, $m=N$; this is the $\gamma$ of Assumption~\ref{bpP:p11a:sel-data}). Further
\[\bar T=1+\frac{16}{\pi s}\Bigl(\lvert A_0\rvert+N\log\frac{32N}s\Bigr),\qquad t_\flat=\min\Bigl\{\frac\vartheta{\sqrt N},\ \frac{\gamma\bar T^{-\tau-1}}{4C_1}\Bigr\},\qquad\bar\Xi=2^{\tau+1}\Bigl(1+\frac1\varsigma+\frac1{t_\flat}+\bar T\varsigma\Bigr),\]
$\theta_{\max}$ as in Definition~\ref{bpP:p11a:adm-params} ($\beta/2$ for $\rho=1$; $\min\{\beta,\beta/(2^{\mu+1}(2M_0+\beta))\}$ for $\rho=2$), and
\[\lambda_\flat=\min\Bigl\{\frac{t_\flat}8,\ \ \min\Bigl\{\frac\varpi2,\frac{2^{-10}}{\bar\Xi}\Bigr\}\frac{\pi\gamma\varsigma^{N+1}}{C_Q2^N\bar T^{\tau}},\ \ \frac{\theta_{\max}t_\flat^\mu}{4c'_{N,\mu}}\Bigr\},\qquad\varepsilon^\flat_{\mathrm{adm}}=\varepsilon^\flat_{\mathrm{adm}}(\mathfrak d):=\frac{\lambda_\flat^2}{2NC_1}.\]
Below $\varepsilon_{\mathrm{adm}}$ without a decoration is the number $\varepsilon_{\mathrm{adm}}^\flat$; the number of Definition~\ref{bpP:p11a:adm-params} is written $\varepsilon_{\mathrm{adm}}^{(3)}$. Both depend on the data and not on $h$ or $f$.
\source{Ours (PLAN-11a.md, rows \texttt{budget}, D-9). $\varepsilon^\flat_{\mathrm{adm}}$ is obtained from $\varepsilon^{(3)}_{\mathrm{adm}}$ by replacing each non-monotone ingredient ($T_0$, $t_0$, $\Xi$, $C_1t_0$) by a monotone bound; it is smaller, so every theorem valid up to $\varepsilon_{\mathrm{adm}}^{(3)}$ is valid up to $\varepsilon^\flat_{\mathrm{adm}}$.}
\end{definition}

The reason for the substitution is that Part~11 needs the threshold as a function of the data that is monotone, so that upper bounds for $M,M_0,C_1,d$ and lower bounds for $\beta,\vartheta,s,\varepsilon^*$ suffice. In $\varepsilon^{(3)}_{\mathrm{adm}}$ the quantity $\Xi$ contains both $t_0/\varsigma_0$ and $\varsigma_0/t_0$, and $C_1t_0$ increases with $C_1$ for small $C_1$.

\begin{proposition}[The budget; monotone minorant and polynomial lower bounds]\label{bpP:p11a:budget}
Let $\mathfrak d$ be as above.
\begin{enumerate}[label=(\alph*)]
\item\label{bpP:p11a:bu-minor} (minorant) $0<\varepsilon_{\mathrm{adm}}^\flat(\mathfrak d)\le\varepsilon_{\mathrm{adm}}^{(3)}(\mathfrak d')$ for every $\mathfrak d'$ that differs from $\mathfrak d$ only by $M_0'\le M_0$.
\item\label{bpP:p11a:bu-mono} (monotone) $\gamma$ and $\varepsilon^\flat_{\mathrm{adm}}$ are nondecreasing in $\beta,\vartheta,s,\varepsilon^*$ and nonincreasing in $M,d$; $\varepsilon^\flat_{\mathrm{adm}}$ is moreover nonincreasing in $M_0$ and in $C_1$ ($\gamma$ does not depend on them). Here $\beta\le M$, $\varepsilon^*<\operatorname{meas}K$, $C_1\ge1$, $\vartheta,s\le1$ are kept, and $d$ is the diameter of $K$.
\item\label{bpP:p11a:bu-poly} (polynomial lower bound, all data) the conclusion of Proposition~\ref{bpP:p11a:adm-budget} holds for $\varepsilon^\flat_{\mathrm{adm}}$ in place of $\varepsilon_{\mathrm{adm}}$, with the same hypotheses, the same exponent $B=2H+a_C$ and a constant $c'$ with the same dependence.
\item\label{bpP:p11a:bu-beta} (the form (K5)) in the situation of Corollary~\ref{bpP:p11a:adm-budget-beta} (all entries but $\beta$ fixed), $\varepsilon^\flat_{\mathrm{adm}}\ge c_*\beta^{B_0}$ for $\beta\in(0,\min\{1,M\}]$, $B_0=2\max\{2G_0,1+\mu G_0\}$, $G_0=(\mu+1)/\rho$ ($B_0=\mu(\mu+1)+2$ for $\rho=2$), $c_*$ depending only on the fixed data.
\end{enumerate}
\source{Ours; File~11a-3, Proposition~\ref{bpP:p11a:adm-budget} and Corollary~\ref{bpP:p11a:adm-budget-beta} for the unmodified threshold.}
\end{proposition}
\begin{proof}
\proofstep{Comparison with File 11a-3}\label{stepP:p11a:qb-comp} Let $T_0,t_0,\Xi,\lambda_{\max}$ be those of Definition~\ref{bpP:p11a:adm-params} for $\mathfrak d'$ (they use $M_0'$ only through $\theta_{\max}$, and $\gamma$ is the same). Since $\max\{1,x\}\le1+\max\{0,x\}$, $T_0\le\bar T$, hence $\gamma T_0^{-\tau-1}\ge\gamma\bar T^{-\tau-1}$ and $t_\flat\le t_0\le1$. Then $C_1t_0/8\ge t_0/8\ge t_\flat/8$ ($C_1\ge1$); $t_0/\varsigma_0\le1/\varsigma$ ($t_0\le1$), $\varsigma_0/t_0\le1/t_\flat$ ($\varsigma\le1$), $T_0\varsigma_0\le\bar T\varsigma$, so $\Xi\le\bar\Xi$ and $\min\{\varpi/2,2^{-10}/\Xi\}\ge\min\{\varpi/2,2^{-10}/\bar\Xi\}$; $T_0^\tau\le\bar T^\tau$ ($\tau\ge0$); $\theta_{\max}(M_0)\le\theta_{\max}(M_0')$ for $M_0'\le M_0$ (the expression $\beta/(2M_0+\beta)$ is decreasing in $M_0$) and $t_\flat^\mu\le t_0^\mu$. Hence each of the three terms of $\lambda_\flat$ is at most the corresponding term of $\lambda_{\max}$, $\lambda_\flat\le\lambda_{\max}$ and $\varepsilon^\flat_{\mathrm{adm}}\le\varepsilon^{(3)}_{\mathrm{adm}}$. Positivity: all ingredients are positive finite numbers.

\proofstep{Monotonicity}\label{stepP:p11a:qb-mono} (i) $\gamma_*$: the first entry increases with $\beta$; $C_{\mathrm{Dio}}=\widehat C\beta^{-1-1/\mu}$ with $\widehat C$ increasing in $M,d$ and decreasing in $\vartheta$ (Lemma~\ref{bpP:p11a:gamma-inverse}), so the second entry increases in $\beta,\varepsilon^*,\vartheta$ and decreases in $M,d$. (ii) $\bar T$ decreases with $s$, and so does $s\bar T=s+\frac{16}\pi(\lvert A_0\rvert+N\log\frac{32N}s)$, whose derivative is $1-16N/(\pi s)<0$ on $(0,1]$. (iii) Hence $\gamma\bar T^{-\tau-1}$ and $t_\flat$ are nondecreasing in $\beta,\vartheta,\varepsilon^*,s$, nonincreasing in $M,d$, and $t_\flat$ is nonincreasing in $C_1$; so $1/t_\flat$, and with $1/\varsigma$ and $\bar T\varsigma=\frac s{24}\bar T$ also $\bar\Xi$, move the opposite way ($\bar\Xi$ is nondecreasing in $C_1$). (iv) In the three terms of $\lambda_\flat$: $t_\flat/8$; $\min\{\cdot\}\,\pi\gamma\varsigma^{N+1}/(C_Q2^N\bar T^\tau)$, with $\varsigma^{N+1}$ increasing and $\bar T^{-\tau}$ increasing in $s$; $\theta_{\max}t_\flat^\mu$, with $\theta_{\max}$ nondecreasing in $\beta$, nonincreasing in $M_0$. So $\lambda_\flat$ has the stated monotonicity and, dividing by $C_1$ in $\lambda_\flat^2/(2NC_1)$, so does $\varepsilon^\flat_{\mathrm{adm}}$.

\proofstep{Polynomial bound, all data}\label{stepP:p11a:qb-poly} Take the hypotheses of Proposition~\ref{bpP:p11a:adm-budget}. Step~\ref{stepP:p11a:ab-gamma} gives $\gamma\ge c_2u^G$. As in Step~\ref{stepP:p11a:ab-T0}: if $a_s=0$ then $s\ge c$ and $\bar T\le C$; if $a_s>0$ then $\lvert A_0\rvert+N\log(32N/s)\le C(1+a_s)u^{-1}$ and $1/s\le c^{-1}u^{-a_s}$; in both cases $\bar T\le Cu^{-L}$, so $\gamma\bar T^{-\tau-1}\ge c_3u^{G+(\tau+1)L}$ and, as in Step~\ref{stepP:p11a:ab-t0}, $t_\flat\ge c_4u^J$ ($J\ge a_\vartheta$ and $J\ge G+(\tau+1)L+a_C$). Then $\bar\Xi\le C(1+u^{-a_s}+u^{-J}+u^{-L})\le Cu^{-X}$ ($1/\varsigma\le Cu^{-a_s}$, $\bar T\varsigma\le\bar T$, $a_s\le L\le X$). The three terms of $\lambda_\flat$ are bounded below exactly as in Step~\ref{stepP:p11a:ab-lam}: $t_\flat/8\ge c_4u^J/8$; $c_8u^{G+(N+1)a_s+\tau L+X}$ (with $\bar T^\tau\le Cu^{-\tau L}$); $c_{10}u^{a_\beta+a_0+\mu J}$ ($\theta_{\max}\ge c_9u^{a_\beta+a_0}$, $t_\flat^\mu\ge c_4^\mu u^{\mu J}$). So $\lambda_\flat\ge c_{11}u^H$ and $\varepsilon^\flat_{\mathrm{adm}}\ge c_{11}^2u^{2H}/(2NC_1)\ge c'u^{2H+a_C}$ by $C_1\le c^{-1}u^{-a_C}$.

\proofstep{The form (K5)}\label{stepP:p11a:qb-beta} Specialise (c) as in Step~\ref{stepP:p11a:abb-spec}: $u=\beta$, $a_\beta=1$, all other rates $0$, $c=\min\{1,\vartheta,s,\varepsilon^*,M^{-1},M_0^{-1},C_1^{-1},d^{-1}\}$.
\end{proof}

\subsection{The case $f\equiv0$}\label{bpP:p11a:sec-q-zero}

Files~11a-3 and 11a-4 are stated for a \emph{positive} bound $\varepsilon$ of $\lvert f\rvert$ and the radius $s_0=s_0(\varepsilon)$ of Lemma~\ref{bpP:p11a:free-radius} degenerates as $\varepsilon\to0$. The case $f\equiv0$ is therefore treated directly; it is needed so that the theorem can be stated with $0\le\sup\lvert f\rvert\le\varepsilon_f$, where the displacement bound of (K6) reads $\Delta=0$ for $\varepsilon_f=0$.

\begin{lemma}[$f\equiv0$]\label{bpP:p11a:zero-case}
Let $\mathfrak d$ and $h$ satisfy the hypotheses of Theorem~\ref{bpP:p11a:weak-kam-quantitative} and $f\equiv0$. Put $H_0=\{b\in K:\ \lvert k\cdot\omega(b)\rvert\ge2\gamma\lvert k\rvert_1^{-\tau}\ \forall k\in\Z^N\setminus\{0\}\}$ and $\mathcal K_b(\theta)=(\theta,b)$. Then $H_0$ is compact, $\operatorname{meas}H_0\ge\operatorname{meas}K-\varepsilon^*$, and for $b\in H_0$: $\mathcal K_b$ is real analytic, primary, an embedded Lagrangian torus map, $\partial_t\mathcal K_b(\theta+tw_b)=X_H(\mathcal K_b(\theta+tw_b))$ with $w_b=\omega(b)$ $(\gamma,\tau)$-Diophantine; the defect is $u_b\equiv0$; and $(b,\theta)\mapsto\mathcal K_b(\theta)$ is injective on $H_0\times[0,1)^N$ with image $H_0\times[0,1)^N$ (in the order (position, action): $[0,1)^N\times H_0$) of measure $\operatorname{meas}H_0$.
\end{lemma}
\begin{proof}
\proofstep{The good set}\label{stepP:p11a:zc-set} Apply Corollary~\ref{bpP:p11a:power-law} with $m=N$, $\tilde\omega_k=\omega$ for all $k$ and $2\gamma$ in place of $\gamma$. Hypothesis \ref{bpP:p11a:D3} is trivial (the difference vanishes); \ref{bpP:p11a:D1}, \ref{bpP:p11a:D2} are (A1), (A2) of the theorem (with $M_0$ not used); and $(\sqrt N\,2\gamma)^\rho\le(\sqrt N\gamma_*)^\rho\le2^{-\rho}\beta/4<\beta/4$. Since $2\gamma=\gamma_*\le\gamma_*$, $H_0$ is compact with $\operatorname{meas}H_0\ge\operatorname{meas}K-\varepsilon^*$.

\proofstep{The tori}\label{stepP:p11a:zc-tori} $\lvert k\cdot\omega(b)\rvert\ge2\gamma\lvert k\rvert_1^{-\tau}\ge\gamma\lvert k\rvert_1^{-\tau}$. With $H=h(I)$, $X_H(\theta,I)=(\partial_IH,-\partial_\theta H)=(\omega(I),0)$, so $\partial_t(\theta+t\omega(b),b)=X_H(\theta+t\omega(b),b)$. $D\mathcal K_b=(\mathrm{id},0)$ is injective and has zero symplectic pullback ($\Omega_0((v,0),(v',0))=0$); $\mathcal K_b(\theta+k)=\mathcal K_b(\theta)+(k,0)$; $\mathcal K_b(\theta)=\mathcal K_b(\theta')$ iff $\theta=\theta'$. The evaluation map is the swap $(b,\theta)\mapsto(\theta,b)$, an isometry of $\R^{2N}$, so it is injective, its image $[0,1)^N\times H_0$ is Borel, of measure $\operatorname{meas}H_0$ \Kref{bp:p1:ext-measure}.
\end{proof}

\subsection{The theorem}\label{bpP:p11a:sec-q-thm}

\begin{theorem}[Quantitative maximal-case analytic KAM theorem under weak non-degeneracy]\label{bpP:p11a:weak-kam-quantitative}
Let $N,\mu\ge1$ be integers, $\rho\in\{1,2\}$, and $\tau=\tau_D>\max\{N\mu/\rho-1,\ N-1\}$. Let $K\subset\R^N$ be compact, $d=\operatorname{diam}_\infty K$, and let $\vartheta,s\in(0,1]$, $C_1\ge1$, $M,M_0>0$, $0<\beta\le\min\{1,M\}$, $\varepsilon^*\in(0,\operatorname{meas}K)$ be given: the \emph{data} $\mathfrak d$. Define $\gamma=\gamma(\mathfrak d)>0$ and $\varepsilon_{\mathrm{adm}}=\varepsilon^\flat_{\mathrm{adm}}(\mathfrak d)>0$ by Definition~\ref{bpP:p11a:thresholds}, so that they depend on $\mathfrak d$ and on nothing else. Then the following holds for every pair $(h,f)$ with:
\begin{enumerate}[label=(A\arabic*)]
\item $U=K+2\vartheta\subset\C^N$ (max norm); $h:U\to\C$ and $f:T_s\times U\to\C$ holomorphic, $f$ $1$-periodic in $\theta$, $h$ and $f$ real on real points; $\omega=\nabla h$ is $C^\infty$ on the real open set $V=U\cap\R^N\supset K_\vartheta$;
\item $\beta_{\rho,\mu}(\omega,K)\ge\beta$ (Definition~\ref{bpP:p11a:index-def});
\item $\lvert D^kg_\rho[\omega](c,y)\rvert_{\mathrm{op}}\le M$ for $c\in S^{N-1}$, $1\le k\le\mu+1$, $y\in K_\vartheta$;
\item $\lvert D^2h\rvert_{\infty\to\infty}\le C_1$ on $U$ and $\lVert\omega\rVert^\mu_{\mathrm{dg},K_\vartheta}\le M_0$,
\end{enumerate}
and every $\varepsilon_f$ with $\sup_{T_s\times U}\lvert f\rvert\le\varepsilon_f\le\varepsilon_{\mathrm{adm}}$ ($\varepsilon_f=0$ allowed), for the Hamiltonian $H(\theta,I)=h(I)+f(\theta,I)$ with $X_H=(\partial_IH,-\partial_\theta H)$ and $\Omega_0=dI\wedge d\theta$:
\begin{enumerate}[label=(K\arabic*)]
\item there is a compact $\mathcal H\subset K$ with $\operatorname{meas}\mathcal H\ge\operatorname{meas}K-\varepsilon^*$;
\item for $b\in\mathcal H$ there is a real analytic $\mathcal K_b:\R^N\to\R^N\times\R^N$ with $\mathcal K_b(\theta+k)=\mathcal K_b(\theta)+(k,0)$ for $k\in\Z^N$ (so $u_b:=\mathcal K_b-(\mathrm{id},b)$ is $\Z^N$-periodic), $D\mathcal K_b$ injective, $\mathcal K_b^*\Omega_0=0$, $\mathcal K_b(\theta)=\mathcal K_b(\theta')$ only if $\theta'-\theta\in\Z^N$, and a vector $w_b\in\R^N$ with $\lvert k\cdot w_b\rvert\ge\gamma\lvert k\rvert_1^{-\tau}$ for all $k\in\Z^N\setminus\{0\}$ and $\partial_t\mathcal K_b(\theta+tw_b)=X_H(\mathcal K_b(\theta+tw_b))$ for all $\theta\in\R^N$, $t\in\R$;
\item $\sup_{b\in\mathcal H,\ \theta\in\R^N}\lvert u_b(\theta)\rvert_\infty\le\Delta(\varepsilon_f;\mathfrak d):=\min\{2^{-18},\ c_\Delta\varepsilon_f^{1/2}/\gamma\}$, with $c_\Delta=c_\Delta(N,\mu,\tau,s,C_1)$ of Theorem~\ref{bpP:p11a:displacement};
\item $\lvert\mathcal K_b(\theta)-\mathcal K_{b'}(\theta')\rvert_\infty\ge(1-2^{-10})\max\{\lvert\theta-\theta'\rvert_\infty,\lvert b-b'\rvert_\infty\}$ for $b,b'\in\mathcal H$, $\theta,\theta'\in\R^N$, and $\mathcal K$ is $(1+2^{-10})$-Lipschitz; so $(b,\theta)\mapsto\mathcal K_b(\theta)$ is injective on $\mathcal H\times[0,1)^N$, its image $\mathcal B$ is Borel and $\operatorname{meas}\mathcal B\ge2^{-2N}\operatorname{meas}\mathcal H$; moreover $\mathcal K(\mathcal H\times[0,1]^N)$ is compact;
\item (budget) $\varepsilon_{\mathrm{adm}}$ is nondecreasing in $\beta,\vartheta,s,\varepsilon^*$, nonincreasing in $M,M_0,C_1,d$, and if all entries of $\mathfrak d$ except $\beta$ are held fixed there is $c_*>0$, depending only on them, with $\varepsilon_{\mathrm{adm}}\ge c_*\beta^{B_0}$ for all $\beta\in(0,\min\{1,M\}]$, $B_0=2\max\{2G_0,1+\mu G_0\}$, $G_0=\frac{\mu+1}\rho$ ($B_0=\mu(\mu+1)+2$ if $\rho=2$) (Proposition~\ref{bpP:p11a:budget}; the general polynomial lower bound in all data is Proposition~\ref{bpP:p11a:budget}\ref{bpP:p11a:bu-poly});
\item (displacement power) with the same entries fixed there is $c_D>0$ depending only on them such that $\sup\lvert u_b\rvert_\infty\le c_D\,\varepsilon_f^{\kappa_D}\beta^{-a_D}$, $\kappa_D=\frac12$, $a_D=G_0$, for all $\beta\in(0,\min\{1,M\}]$.
\end{enumerate}
The map $\mathcal K_b$ is primary, not periodic: periodic maps appear only after the charts of Parts~11 and 13. The Diophantine property of $w_b$, the Lipschitz bound of (K4) and the polynomial threshold of (K5) are conclusions.
\source{Rus01 Thm 1.7 (p.~126), \S14.10 (p.~171) and \S\S15--18 (the scheme) as modified in Files~11a-1 to 11a-4 (D-1 to D-22); CP09 Thm 3, Lemma 8 and (34)--(47) (pp.~6--7, 18--21) for the shape of the statement. The displacement, the Lipschitz family and the explicit $\varepsilon_{\mathrm{adm}}$ are not in Rus01 Thm 1.7 (A2--A4 of the Russmann check).}
\end{theorem}
\begin{proof}
\proofstep{Totality and quantifier order}\label{stepP:p11a:qt-order} $\gamma$ and $\varepsilon_{\mathrm{adm}}$ are explicit positive numbers of $\mathfrak d$ (Definition~\ref{bpP:p11a:thresholds}); $\varepsilon^*<\operatorname{meas}K$ forces $\operatorname{meas}K>0$ and $d>0$. Nothing in (A1)--(A4) or in $\varepsilon_f\le\varepsilon_{\mathrm{adm}}$ mentions a conclusion: no Diophantine property, no frequency $w_b$, no torus, no Lipschitz bound. The Euclidean thickening $K_\vartheta$ is contained in $V$ because Euclidean distance $\le\vartheta$ implies max-distance $\le\vartheta<2\vartheta$.

\proofstep{Reduction to the standing data of File 11a-3}\label{stepP:p11a:qt-reduce} Let $M_0^{\mathrm{act}}:=\lVert\omega\rVert^\mu_{\mathrm{dg},K_\vartheta}\le M_0$ and $\mathfrak d'$ the datum with $M_0$ replaced by $M_0^{\mathrm{act}}$ (if $M_0^{\mathrm{act}}=0$ replace it by any positive number $\le M_0$; $M_0$ enters only through the upper bound $\theta_{\max}$). Then (A1)--(A4) are Assumption~\ref{bpP:p11a:sel-data}(c),(d) for $(\mathfrak d',h,f)$ with the $\varepsilon$ of that assumption any positive number with $\sup\lvert f\rvert\le\varepsilon$, and $\gamma(\mathfrak d')=\gamma(\mathfrak d)$ ($\gamma$ does not depend on $M_0$). By Proposition~\ref{bpP:p11a:budget}\ref{bpP:p11a:bu-minor}, $\varepsilon_{\mathrm{adm}}^\flat(\mathfrak d)\le\varepsilon^{(3)}_{\mathrm{adm}}(\mathfrak d')$.

\proofstep{The case $\varepsilon_f=0$}\label{stepP:p11a:qt-zero} Then $f\equiv0$ and Lemma~\ref{bpP:p11a:zero-case} gives (K1), (K2) with $\mathcal H=H_0$, (K3) (with $u_b=0$), (K4) (the evaluation map is the swap, $\lambda=0$, $\operatorname{meas}\mathcal B=\operatorname{meas}\mathcal H$), and (K6) trivially ($0\le0$). Assume now $\varepsilon_f>0$.

\proofstep{The iteration}\label{stepP:p11a:qt-iter} Take $\varepsilon:=\varepsilon_f\le\varepsilon_{\mathrm{adm}}\le\varepsilon^{(3)}_{\mathrm{adm}}(\mathfrak d')$ and the situation $(\mathrm S)$ of Assumption~\ref{bpP:p11a:situation} for $\mathfrak d'$, i.e.\ the schedule of Definition~\ref{bpP:p11a:schedule} with the radius $s_0(\varepsilon)$ of Lemma~\ref{bpP:p11a:free-radius} (this radius is chosen here, from $\varepsilon$ and the data; it is not an input of the theorem). Theorem~\ref{bpP:p11a:parameter-selection} runs the iteration for all stages and gives the compact set $H_\infty\subset K_\infty\subset K$ with $\operatorname{meas}H_\infty\ge\operatorname{meas}K-\varepsilon^*$, and $\omega_*(b)$ $(\gamma,\tau)$-Diophantine for $b\in H_\infty$. Put $\mathcal H=H_\infty$, $w_b=\omega_*(b)$. This is (K1).

\proofstep{(K2)}\label{stepP:p11a:qt-k2} Theorem~\ref{bpP:p11a:limit-torus}(a)--(e), for $b\in K_\infty\supset\mathcal H$: $\mathcal K_b$ real analytic, primary, $D\mathcal K_b w_b=X_H(\mathcal K_b)$ (hence $\partial_t\mathcal K_b(\theta+tw_b)=X_H(\mathcal K_b(\theta+tw_b))$), Lagrangian, $D\mathcal K_b$ injective, $\mathcal K_b(\theta)=\mathcal K_b(\theta')$ only if $\theta=\theta'$ (a fortiori mod $\Z^N$), and on $H_\infty$ the Diophantine bound.

\proofstep{(K3)}\label{stepP:p11a:qt-k3} Theorem~\ref{bpP:p11a:displacement} gives $\sup\lvert u_b\rvert_\infty\le\Delta(\varepsilon)\le\min\{2^{-18},c_\Delta\varepsilon^{1/2}/\gamma\}$; $c_\Delta$ depends on $(N,\mu,\tau,s,C_1)$.

\proofstep{(K4)}\label{stepP:p11a:qt-k4} Theorem~\ref{bpP:p11a:evaluation-bilipschitz}(a),(c),(d) with $\lambda_*=2^{-10}$, $H=H_\infty$: the bi-Lipschitz bounds (in the pair $(\theta,b)$ with the max norm of the two blocks; (K4) states the lower bound for the max of the two differences, which is the same quantity), injectivity on $\mathcal H\times\R^N$, Borel image with $\operatorname{meas}\mathcal B\ge(1-2^{-10})^{2N}\operatorname{meas}\mathcal H\ge2^{-2N}\operatorname{meas}\mathcal H$, compactness.

\proofstep{(K5)}\label{stepP:p11a:qt-k5} Proposition~\ref{bpP:p11a:budget}\ref{bpP:p11a:bu-mono}, \ref{bpP:p11a:bu-beta}.

\proofstep{(K6)}\label{stepP:p11a:qt-k6} Corollary~\ref{bpP:p11a:displacement-power}: $\sup\lvert u_b\rvert_\infty\le c_\Delta c_\gamma\varepsilon^{1/2}\beta^{-G_0}$ for all $\beta\in(0,\min\{1,M\}]$ and $0<\varepsilon\le\varepsilon_{\mathrm{adm}}$, with $c_D=c_\Delta c_\gamma$ depending on the fixed entries only (and $c_\Delta$ on $C_1,s$, which are among them). It applies with $\varepsilon=\varepsilon_f$ since $\varepsilon_{\mathrm{adm}}\le\varepsilon^{(3)}_{\mathrm{adm}}$, the threshold of File~11a-4.
\end{proof}

\paragraph{Why every constant is chosen in the right order.} $\gamma$ is chosen from $(\mathfrak d)$ through $\gamma_*$ (File~11a-1) before the iteration, because the measure loss $C_{\mathrm{Dio}}\gamma^{\rho/\mu}$ must be $\le\varepsilon^*$ for the mode-wise good set (Theorem~\ref{bpP:p11a:parameter-selection}(b)). $\varepsilon_{\mathrm{adm}}$ depends on $\gamma$ because the collar $t_0$ and the cutoff $T_0$ are tied to the divisor $\gamma T_0^{-\tau-1}$; the initial action radius $s_0\asymp\varepsilon^{1/2}$ is then chosen from $\varepsilon$ (File~11a-4), which makes $\Delta\asymp\varepsilon^{1/2}/\gamma$ and gives (K6).

\subsection{The form consumed by Parts 11 and 13}\label{bpP:p11a:sec-q-phys}

\begin{corollary}[Physical interface]\label{bpP:p11a:physical-interface}
In the situation of Theorem~\ref{bpP:p11a:weak-kam-quantitative} let $S:=\mathcal K(\mathcal H\times[0,1)^N)=\{\mathcal K_b(\theta):b\in\mathcal H,\theta\in[0,1)^N\}\subset\R^N\times\R^N$ (positions, actions) and $\widetilde S:=\mathcal K(\mathcal H\times\R^N)=S+(\Z^N\times\{0\})$. Then:
\begin{enumerate}[label=(\alph*)]
\item $S$ is Borel and $0<2^{-2N}(\operatorname{meas}K-\varepsilon^*)\le2^{-2N}\operatorname{meas}\mathcal H\le\operatorname{meas}S\le(1+2^{-10})^{2N}\operatorname{meas}\mathcal H<\infty$;
\item every point of $S$ is within max-distance $\Delta=\Delta(\varepsilon_f;\mathfrak d)\le2^{-18}$ of $(\theta,b)$ for some $b\in\mathcal H$, $\theta\in[0,1)^N$; in particular its action coordinate lies in $\{I:\operatorname{dist}_\infty(I,\mathcal H)\le\Delta\}\subset K+2^{-18}$;
\item $\widetilde S$ is invariant under the flow of $X_H$: for $z=\mathcal K_b(\theta)\in\widetilde S$ the integral curve of $X_H$ through $z$ exists for all $t\in\R$ and equals $t\mapsto\mathcal K_b(\theta+tw_b)\in\widetilde S$; the quotient of $\widetilde S$ by $\Z^N\times\{0\}$ is compact if $\mathcal H$ is (it is);
\item $\widetilde S$ is the disjoint union over $b\in\mathcal H$ of the embedded Lagrangian tori $\mathcal K_b(\R^N)/\!\sim$, each with the Diophantine frequency $w_b$ of (K2).
\end{enumerate}
\source{Ours (PLAN-11a.md, row \texttt{physical-interface}); period-one form of the frozen torus clauses minus the chart-specific ones (those are verified in Part~13 after the polar and Birkhoff charts, Theorem~\ref{bpP:p1:torus-transport}).}
\end{corollary}
\begin{proof}
\proofstep{(a)}\label{stepP:p11a:pi-a} Theorem~\ref{bpP:p11a:evaluation-bilipschitz}(c) for the lower and upper bounds; $\operatorname{meas}\mathcal H\le\operatorname{meas}K<\infty$ (compact $K$) gives finiteness, and $\operatorname{meas}\mathcal H\ge\operatorname{meas}K-\varepsilon^*>0$ by (K1).

\proofstep{(b)}\label{stepP:p11a:pi-b} $\mathcal K_b(\theta)=(\theta,b)+u_b(\theta)$ and $\lvert u_b\rvert_\infty\le\Delta\le2^{-18}$ (K3). (The inclusion in $K+2^{-18}$ is only a statement about max-distance; it is not claimed that $2^{-18}<\vartheta$.)

\proofstep{(c)}\label{stepP:p11a:pi-c} By (K2), $t\mapsto\mathcal K_b(\theta+tw_b)$ is an integral curve of $X_H$ defined on $\R$ (its action coordinates differ from $b$ by less than $s_0\le\vartheta$, Theorem~\ref{bpP:p11a:displacement} and Lemma~\ref{bpP:p11a:free-radius}, so the curve lies in the domain of $H$, as used in Theorem~\ref{bpP:p11a:limit-torus}(b)), and $X_H$ is real analytic hence locally Lipschitz on the real part of $T_s\times U$, so by uniqueness (Picard--Lindel\"of, \Kref{bp:p1:ext-mvt} for the Lipschitz bound) it is the integral curve through $z$. The curve stays in $\widetilde S$ because $\mathcal K_b(\theta+tw_b)\in\mathcal K(\mathcal H\times\R^N)$. Compactness of the quotient is Theorem~\ref{bpP:p11a:evaluation-bilipschitz}(d).

\proofstep{(d)}\label{stepP:p11a:pi-d} The tori for different $b$ are disjoint by injectivity of $\mathcal K$ on $\mathcal H\times\R^N$ (K4); each is embedded mod $\Z^N$ with the properties (K2).
\end{proof}

\subsection{Sources, deviations, check of the interface, what is received}\label{bpP:p11a:notes5}

\paragraph{Sources read.} Rus01 Thm 1.7 (p.~126) and \S14.10 (p.~171) as recorded in PLAN-11a.md; CP09 Thm 3 and Lemma 8 as recorded there (the printed constants are not reproduced; the explicit thresholds are Definition~\ref{bpP:p11a:thresholds}). Files~11a-1 (Corollary~\ref{bpP:p11a:power-law}), 11a-3 (Definition~\ref{bpP:p11a:adm-params}, Theorem~\ref{bpP:p11a:parameter-selection}, Proposition~\ref{bpP:p11a:adm-budget}, Corollary~\ref{bpP:p11a:adm-budget-beta}), 11a-4 (all of Section ``What File 11a-5 must take'', items (U1)--(U4)), Part~11 Definition~\ref{bpP:p11:kam-input}, Part~13 (forward references). I did not re-read Rus01 pages in this file; nothing here depends on them beyond the cited Files.

\paragraph{Check of Definition~\ref{bpP:p11:kam-input} item by item.}
\begin{itemize}
\item \emph{Datum.} $\mathfrak d_{\mathrm{P11}}=(\mu,\beta,K,\vartheta,s,M,M_0,C_1,\tau_D,\varepsilon^*)$ is the datum $\mathfrak d$ of Theorem~\ref{bpP:p11a:weak-kam-quantitative} with $N=2n$, $\rho=2$; the requirements $\vartheta,s\in(0,1]$, $C_1\ge1$, $M_0>0$, $\varepsilon^*<\operatorname{meas}K$, $0<\beta\le\min\{1,M\}$, $\tau_D>\max\{N\mu/2-1,N-1\}$ coincide, and $M\ge\beta$ follows from $\beta\le\min\{1,M\}$.
\item \emph{(A1)--(A4)} are (A1)--(A4) of the theorem verbatim. The only difference is that Part~11 writes $\lVert\omega\rVert^\mu_{\mathrm{dg},K_\vartheta}\le M_0$ and File~11a-3 uses the actual value; Step~\ref{stepP:p11a:qt-reduce} bridges this by the monotonicity in $M_0$.
\item \emph{(K1)--(K4)}: Steps~\ref{stepP:p11a:qt-iter}--\ref{stepP:p11a:qt-k4}; $\Delta(\varepsilon_f;\mathfrak d)$ is the number of (K3). \emph{(K5)}: the theorem's item (K5); $c_*$ depends only on the fixed entries. \emph{(K6)}: item (K6) with $\kappa_D=\frac12$, $a_D=\frac{\mu+1}2$, $c_D$ as stated.
\item \emph{Mismatches.} (1) Part~11 says ``$\sup\lvert f\rvert\le\varepsilon_f\le\varepsilon_{\mathrm{adm}}$'' without $\varepsilon_f>0$; the theorem covers $\varepsilon_f=0$ by Lemma~\ref{bpP:p11a:zero-case}, so nothing is needed. (2) Part~11 says $\varepsilon_{\mathrm{adm}}\ge c_*\beta^{B_0}$ ``proved by Corollary~\ref{bpP:p11a:adm-budget-beta}'': that corollary is for $\varepsilon^{(3)}_{\mathrm{adm}}$; the threshold of the theorem is $\varepsilon^\flat_{\mathrm{adm}}\le\varepsilon^{(3)}_{\mathrm{adm}}$, for which the same bound is Proposition~\ref{bpP:p11a:budget}\ref{bpP:p11a:bu-beta}. The citation in Part~11 should be to \texttt{bpP:p11a:budget}; the value $B_0$ is unchanged. (3) The data of Part~11 include $d=\operatorname{diam}_\infty K$ and $M_0$ as fixed entries; they are (Remark: $d=\operatorname{diam}_\infty[-\tfrac12,\tfrac12]^{2n}=1$).
\end{itemize}

\paragraph{Check of the forward references of Part 13.} Part~13 cites \texttt{weak-kam-quantitative}, \texttt{limit-torus}, \texttt{displacement}, \texttt{evaluation-bilipschitz}, \texttt{lipschitz-family}, \texttt{budget}, \texttt{physical-interface}: all exist (the first, sixth, seventh here; the others in File~11a-4). \texttt{reference-constants} of the plan is not written: Part~11 does not use it (Lemma~\ref{bpP:p11:amount}) and Part~13 does not cite it.

\paragraph{Deviations.} \textbf{D-23}: the threshold of the theorem is the monotone minorant $\varepsilon^\flat_{\mathrm{adm}}$ of File~11a-3's $\varepsilon_{\mathrm{adm}}$, not the latter, so that the monotonicity required by the plan holds; the polynomial exponent $B$ is unchanged. \textbf{D-24}: the case $f\equiv0$ is treated separately (Lemma~\ref{bpP:p11a:zero-case}) because File~11a-4 is stated for a positive bound; $\varepsilon_f$ plays the role of $\varepsilon$ otherwise. \textbf{D-25}: the second lower bound $\tau_D>N-1$ (the frozen clause (vii) $\tau>d-1$) is part of the data of the theorem but is not used in any proof of Part~11a.

\paragraph{Suspected errata.} None in this file. (Those of Files~11a-1 to 11a-4 are recorded there.)

\paragraph{Doubts about frozen statements.} None: this part contains no frozen statement; the frozen torus clauses are assembled in Part~13.

\paragraph{External facts used.} \Kref{bp:p1:ext-measure}, \Kref{bp:p1:ext-mvt}; elementary calculus (derivative sign of $s+c\log(1/s)$ in Step~\ref{stepP:p11a:qb-mono}); Picard--Lindel\"of uniqueness for locally Lipschitz fields (Corollary~\ref{bpP:p11a:physical-interface}(c), listed here once); the other external facts are those of Files~11a-1 to 11a-4.

\paragraph{Forward references.} None. All labels cited exist in Parts~11a-1 to 11a-4; Parts~11 and 13 cite this file.

\paragraph{Label index.} \texttt{thresholds}, \texttt{budget}, \texttt{zero-case}, \texttt{weak-kam-quantitative}, \texttt{physical-interface}.
\par\endgroup
